
\subsection{Distance Assignment}\label{sec:distance-treatment}

Implementing the distance treatment required jointly selecting treatment locations and targeted communities under geographic and logistical constraints. We relied on geospatial data on primary schools, churches, and mosques as spatial reference points: areas around these locations provide feasible centralized locations for points of treatment (PoTs), and the dense geographic coverage of primary schools and religious centers in the study area provides a proxy for locating nearby communities, in the absence of geo-coded community locations.

We proceed in three steps: (i) select PoT-centered clusters using a geographically constrained procedure, (ii) randomly assign each selected cluster to a Close or Far distance condition, (iii) randomly select one targeted community per cluster whose centroid lies in the assigned distance band relative to that cluster's PoT. Appendix~\ref{sec:site-selection} provides full details. %algorithmic 

\vspace{0.4em}
\noindent\textbf{1.\ Cluster (PoT) selection.} We begin from a sampling frame of 1{,}451 potential clusters, each defined by a 2.5~km radius catchment circle centered on a primary school/religious center (the candidate PoT). We select 158 clusters using a sequential acceptance--rejection algorithm: we repeatedly draw a candidate cluster at random and accept if, after accounting for overlap with previously selected clusters, all selected clusters retain a sufficiently large non-overlapping catchment.\footnote{We targeted 150 clusters and drew eight additional clusters as fallbacks. Field feasibility constraints limited implementation to 144 clusters. Eight clusters were dropped during mapping because the anchor school/religious center or its intended location could not be verified (this affected, two clusters in Busia ((far; bracelet/calendar), (close; ink/control)), two clusters in Siaya ((far; ink/control), (far; bracelet/calendar)) and four clusters in Kakamega (3 (far; bracelet/calendar), 1 (close; ink/control)) and six of the original 150 clusters were dropped because implementation was practically not feasible as CHVs were not available, the local chief was against, or community was adversarial towards the program: two clusters in Siaya (Far; ink and far; Ink/Control), two clusters in Busia (Close; Calendar, and Close; Control) and 2 cluster in Kakamega (Far; Calendar and Far; Control).} %TIGHTEN AND CHECK BALANCE -- WHAT IS IT THAT WE NEED TO SHOW? BASELINE BALANCE ACROSS 144 should be okay; however, could be useful to justify WHY we have so many fewer FAR clusters. We need the original assignment for that. 

We operationalize this feasibility requirement by counting available primary schools/religious centers in the non-overlapping portion of each 2.5~km catchment: a location counts as available if at least 60\% of a 1~km radius circle around it lies within the non-overlapping portion of the catchment. We require each selected cluster to retain at least one such point. When assessing overlap, we use buffer circles larger than 2.5~km (3--4~km) to provide separation between nearby clusters. This implies that communities selected from the non-overlapping catchment are at least 0.5~km closer to their own assigned PoT than to the nearest alternative selected PoT. Figure~\ref{fig:sites} maps the selected PoTs and their non-overlapping catchment areas.

\vspace{0.4em}
\noindent\textbf{2.\ Close/Far assignment.}
After selecting clusters, we randomly assign them, stratified by county, to a Close or Far distance condition. The distance category is defined in terms of the distance from the centroid of the ultimately selected community to the cluster's PoT: 0--1.25~km in Close clusters and 1.25--2.5~km in Far clusters. The 1.25~km cutoff was chosen ex ante as the midpoint of the 2.5~km catchment radius and corresponds to roughly half the upper range of distances adults were generally willing to travel for deworming treatment in our setting.
%MISSING: explain why 80 Close and 64 Far clusters in the end. "Footnote (recommended): “Clusters were randomly assigned to Close or Far within county. Field feasibility constraints limited implementation to 144 of the 158 selected clusters, so exact 1:1 balance in the realized sample was not imposed. The final sample contains 80 Close and 64 Far clusters (two‑sided binomial p=0.21)" -- we can show balance. But did we drop clusters differentially? Where is that code? Email Karim and Ed about this!

\vspace{0.4em}
\noindent\textbf{3.\ Targeted community selection.}
To implement the assigned distance category, we selected for each cluster, from its non-overlapping area and according to its distance assignment, a primary school/religious center as an anchor to locate eligible communities to target. In Close clusters, the anchor is typically the PoT location (or a nearby school/religious center). In Far clusters, the anchor is chosen in the outer ring so that listed communities lie 1.25--2.5~km from the PoT. Enumerators listed all eligible communities near the anchor location and we randomly selected one targeted community per cluster. In all but two clusters, the assigned PoT was the closest deworming site to the community. 